\documentclass[11pt,a4paper]{article}
\usepackage[english]{cursus}
\author{dr. Maarten Slembrouck}

\begin{document}

\labtitle{Lab 5 -- Data Visualization I}{Tell the story with plotly}{Course text for the notebook \code{lab5\_data\_visualization\_I.py}}

\section*{What is this lab about?}

In two labs you turned messy exports into tidy, combined tables. But a
coach does not read a table with 111 rows. A good chart shows at a
glance what is hidden in the table: who is fastest, who is improving,
how hard a training session was.

\textbf{After this lab you can} choose a chart type that fits a
question, build the standard charts with \code{plotly.express}, show
extra information with colour, facets and size, and make charts
interactive in marimo.

All charts in this text were made with the code above them, on the
club's data. In the notebook they are interactive: you can zoom in and
hover over a point for more information.

\begin{python}[hide]
import pandas as pd
import plotly.express as px
import plotly.io as pio

pio.templates.default = "plotly_white"

def clean_positions(position_column):
    position_map = {
        "forward": "Forward", "fwd": "Forward",
        "midfielder": "Midfielder", "mid": "Midfielder",
        "defender": "Defender", "def": "Defender",
        "goalkeeper": "Goalkeeper", "gk": "Goalkeeper",
    }
    return position_column.str.strip().str.lower().map(position_map)

def fix_sprint_units(seconds_column):
    return seconds_column.where(seconds_column < 20, seconds_column / 1000)

REFERENCE_DATE = pd.Timestamp("2026-09-14")
athletes = pd.read_csv(f"{DATA_DIR}/athletes.csv").assign(
    position=lambda d: clean_positions(d["position"]),
    weight_kg=lambda d: d["weight_kg"].str.replace(",", ".").astype(float),
    age=lambda d: (REFERENCE_DATE - pd.to_datetime(d["birthdate"])).dt.days // 365,
).drop_duplicates(subset=["name", "birthdate"], keep="first")
tests = pd.read_csv(f"{DATA_DIR}/tests.csv").assign(
    sprint_10m_s=lambda d: fix_sprint_units(d["sprint_10m_s"]),
    sprint_30m_s=lambda d: fix_sprint_units(d["sprint_30m_s"]),
)
combined = athletes.merge(tests, on="athlete_id", how="inner")
hr_session = pd.read_csv(f"{DATA_DIR}/hr_session.csv")
hr_session["timestamp"] = pd.to_datetime(hr_session["date"]) + pd.to_timedelta(hr_session["time_s"], unit="s")
hr_with_age = hr_session.merge(athletes[["athlete_id", "age", "name"]], on="athlete_id", how="left")
hr_with_age["hr_max"] = 220 - hr_with_age["age"]
hr_with_age["pct_hr_max"] = 100 * hr_with_age["heart_rate_bpm"] / hr_with_age["hr_max"]
hr_zoned = hr_with_age.assign(hr_zone=lambda d: pd.cut(
    d["pct_hr_max"], bins=[0, 60, 70, 80, 90, 1000],
    labels=["Zone 1", "Zone 2", "Zone 3", "Zone 4", "Zone 5"]))
\end{python}

% ---------------------------------------------------------------------------
\section{The data for this lab}

The notebook starts with one cell that repeats the complete cleaning of
labs 3 and 4. After that you have two tables:

\begin{itemize}
  \item \code{combined}: one row per athlete per test moment, with
        position, gender, weight, sprint times, jump height and Yo-Yo
        level;
  \item \code{hr\_zoned}: the training session per second, with heart
        rate, speed, name and heart-rate zone.
\end{itemize}

\begin{python}
combined.shape, hr_zoned.shape
\end{python}

Both are \textbf{long}: one row per measurement. That is exactly the
shape \code{plotly.express} expects (see lab 4).

The notebook also fixes one colour per position up front, in a
dictionary. That way ``Forward'' has the same colour in every chart, and
the reader only has to learn the legend once.

\begin{python}[noout]
POSITION_COLORS = {
    "Forward": "#2a78d6",
    "Midfielder": "#eb6834",
    "Defender": "#1baf7a",
    "Goalkeeper": "#eda100",
}
\end{python}

You write a colour as a \emph{hex code}: a hash sign followed by six
characters, two for red, two for green and two for blue.

% ---------------------------------------------------------------------------
\section{Principles: the right chart for the right question}

Every chart answers a question. So start from the question, not from the
chart type you happen to know.

\begin{center}
\rowcolors{2}{tint}{white}
\begin{tabularx}{\textwidth}{@{}L{2.8cm}Y L{3.3cm}@{}}
\toprule
\textbf{Kind of question} & \textbf{Example} & \textbf{Suitable chart} \\
\midrule
Distribution & How are sprint times spread out? & histogram, box plot \\
Comparison & Which position jumps highest on average? & bar chart, box plot \\
Relationship & Does a faster sprint mean a higher jump? & scatter plot \\
Evolution & Is this athlete's heart rate climbing? & line chart \\
\bottomrule
\end{tabularx}
\end{center}

\subsection{Honest axes}

A chart can mislead as easily as it informs. The best-known trap is a
\textbf{truncated axis}. Below are the same means twice. In the first
chart the y-axis starts at 0:

\begin{python}[fig=honest]
sprint_means = combined.groupby("position", as_index=False)["sprint_30m_s"].mean()
fig = px.bar(
    sprint_means, x="position", y="sprint_30m_s",
    color="position", color_discrete_map=POSITION_COLORS,
    title="Average 30 m sprint time by position",
    labels={"sprint_30m_s": "Average sprint time (s)", "position": "Position"},
)
fig.update_yaxes(rangemode="tozero")
\end{python}

\code{as\_index=False} keeps \code{position} as an ordinary column
(instead of making it the index), so you can use it as \code{x}. In the
second chart the y-axis starts just below the smallest value:

\begin{python}[fig=misleading]
fig = px.bar(
    sprint_means, x="position", y="sprint_30m_s",
    color="position", color_discrete_map=POSITION_COLORS,
    title="MISLEADING - average 30 m sprint time by position",
    labels={"sprint_30m_s": "Average sprint time (s)", "position": "Position"},
)
fig.update_yaxes(range=[5.1, 5.55])
\end{python}

The same small differences (at most 0.3 seconds) now look huge. A bar
shows a size with its \emph{length}; so that length has to start at
zero. For the same reason you never use two different y-axes in one
chart: the reader cannot compare the two scales.

% ---------------------------------------------------------------------------
\section{Making charts with \code{plotly.express}}

\code{plotly.express} is a library for making interactive charts with
little code. You import it as \code{px}:

\begin{python}[noout]
import plotly.express as px
\end{python}

Every function follows the same pattern:

\begin{python}[noeval]
fig = px.KIND(table, x="column", y="column", color="column",
              title="...", labels={"column": "Readable name"})
\end{python}

\begin{itemize}
  \item the \textbf{table} comes first; after that you refer to columns
        by their name, as text;
  \item \code{title} is the title above the chart;
  \item \code{labels} is a dictionary that translates column names into
        readable axis labels, with a unit: ``Sprint time (s)'' instead of
        \code{sprint\_30m\_s};
  \item the result is a \textbf{figure} object. Put the variable on the
        last line of a cell and marimo shows the chart.
\end{itemize}

\subsection{Histogram: how is one variable distributed?}

A histogram divides the values into bins and shows how many values fall
in each bin. \code{nbins} chooses the (approximate) number of bins.

\begin{python}[fig=histogram]
fig = px.histogram(
    athletes, x="height_cm", nbins=15,
    title="Distribution of athlete height",
    labels={"height_cm": "Height (cm)"},
)
\end{python}

A histogram needs only one column (\code{x}); the y-axis is the number
of athletes per bin, which plotly counts itself. Play with
\code{nbins}: too few bins hide detail, too many make the chart messy.

\subsection{Box plot: comparing distributions}

A box plot summarizes a distribution in five numbers. This is how you
read it:

\begin{itemize}
  \item the \textbf{line in the box} is the median: half of the values
        are below it, half above;
  \item the \textbf{box} runs from the first to the third quartile: the
        middle half of the values;
  \item the \textbf{whiskers} run to the lowest and highest value that
        is not exceptionally far from the box;
  \item \textbf{separate points} are outliers.
\end{itemize}

\begin{python}[fig=boxplot]
fig = px.box(
    combined, x="position", y="sprint_30m_s",
    color="position", color_discrete_map=POSITION_COLORS,
    title="30 m sprint time by position",
    labels={"sprint_30m_s": "Sprint time (s)", "position": "Position"},
    hover_data=["name"],
)
\end{python}

\code{hover\_data=["name"]} adds the name to the information box that
appears when you hover over a point. In the notebook you immediately see
who that outlier is.

\subsection{Scatter plot: is there a relationship?}

In a scatter plot every row is a point, with one variable on the x-axis
and the other on the y-axis.

\begin{python}[fig=scatter]
fig = px.scatter(
    combined, x="sprint_30m_s", y="cmj_height_cm",
    color="position", color_discrete_map=POSITION_COLORS,
    title="Sprint time vs. jump height",
    labels={"sprint_30m_s": "30 m sprint time (s)", "cmj_height_cm": "Jump height (cm)"},
    hover_data=["name", "test_moment"],
)
\end{python}

The points form a broad cloud without a clear direction. If there were a
strong relationship, they would lie around a sloping line: going down if
faster sprinters (a lower time) also jump higher. That is not the case
here: in this data the sprint time says little about the jump height.
``No clear relationship'' is also an answer to the coach's question. And
when you do see a relationship, remember that it is not yet proof that
one causes the other.

\subsection{Line chart: how does something change?}

A line chart connects consecutive points. For that, the data has to be
\textbf{sorted} on the x-axis, otherwise the line jumps back and forth.
With \code{color="athlete\_id"} every athlete gets their own line:

\begin{python}[fig=line]
fig = px.line(
    combined.sort_values("test_moment"), x="test_moment", y="sprint_30m_s",
    color="athlete_id",
    title="30 m sprint time across the season, per athlete",
    labels={"sprint_30m_s": "Sprint time (s)", "test_moment": "Test moment"},
)
fig.update_layout(showlegend=False)
\end{python}

\code{update\_layout(showlegend=False)} hides the legend: with 42
athletes it would take up more space than the chart itself. This works
without \code{pivot\_table}, because \code{combined} is already long.

% ---------------------------------------------------------------------------
\section{More dimensions: colour, facets and size}

A scatter plot shows two variables. With three extra arguments you can
show more:

\begin{itemize}
  \item \code{color="column"}: a category as colour;
  \item \code{facet\_col="column"}: one small chart per value, side by
        side, with the same axes;
  \item \code{size="column"}: a number as the size of the point.
\end{itemize}

\begin{python}[fig=facet]
fig = px.scatter(
    combined, x="sprint_30m_s", y="cmj_height_cm",
    color="position", color_discrete_map=POSITION_COLORS,
    facet_col="gender", size="weight_kg",
    title="Sprint vs. jump height, by position and gender",
    labels={"sprint_30m_s": "30 m sprint time (s)", "cmj_height_cm": "Jump height (cm)"},
)
\end{python}

Five variables in one chart: sprint time, jump height, position, gender
and weight. Notice the third, empty panel \code{gender=nan}: a few rows
have no gender, and plotly makes a separate panel for them. If you don't
want that, drop those rows first with
\code{combined.dropna(subset=["gender"])}, and rename the panel titles
with \code{labels=\{"gender": "Gender"\}}. Facets are often clearer than
cramming everything into one chart, because you can compare the panels
side by side. But more is not always better: only add a dimension if it
helps answer the question.

% ---------------------------------------------------------------------------
\section{Heart rate as a time series, with zones}

A line chart of the heart rate becomes much easier to read when you show
the heart-rate zones as coloured bands in the background. That is what
\code{add\_hrect} is for: a horizontal rectangle between two values on
the y-axis.

\begin{python}[fig=heartrate]
ZONE_COLORS = {"Zone 1": "#cde2fb", "Zone 2": "#86b6ef", "Zone 3": "#3987e5",
               "Zone 4": "#1c5cab", "Zone 5": "#0d366b"}
one_athlete_hr = hr_zoned.loc[hr_zoned["athlete_id"] == "A001"].sort_values("timestamp")
hr_max = one_athlete_hr["hr_max"].iloc[0]

fig = px.line(
    one_athlete_hr, x="timestamp", y="heart_rate_bpm",
    title=f"Heart rate during training - {one_athlete_hr['name'].iloc[0]}",
    labels={"heart_rate_bpm": "Heart rate (bpm)", "timestamp": "Time"},
)
zone_edges_pct = [0, 60, 70, 80, 90, 100]
zone_names = ["Zone 1", "Zone 2", "Zone 3", "Zone 4", "Zone 5"]
for lo, hi, zone in zip(zone_edges_pct[:-1], zone_edges_pct[1:], zone_names):
    fig.add_hrect(y0=lo / 100 * hr_max, y1=hi / 100 * hr_max,
                  fillcolor=ZONE_COLORS[zone], opacity=0.35, line_width=0,
                  annotation_text=zone, annotation_position="right")
fig
\end{python}

The loop deserves some explanation:

\begin{itemize}
  \item \code{zone\_edges\_pct[:-1]} is the list without its last
        element (0, 60, 70, 80, 90): the lower edges.
        \code{zone\_edges\_pct[1:]} is the list without its first element
        (60, 70, 80, 90, 100): the upper edges.
  \item \code{zip(...)} walks through three lists at once and gives one
        element of each list every time: first \code{(0, 60, "Zone 1")},
        then \code{(60, 70, "Zone 2")}, and so on.
  \item \code{lo / 100 * hr\_max} turns a percentage into a heart rate in
        beats per minute.
\end{itemize}

In the notebook you use \code{\_lo} instead of \code{lo}: in marimo a
name that starts with \code{\_} stays inside the cell, so it does not
clash with variables in other cells.

% ---------------------------------------------------------------------------
\section{Interactivity in marimo}

\subsection{UI elements and \code{.value}}

marimo has \textbf{UI elements}: dropdowns, sliders, text boxes. A UI
element is an ordinary Python variable. Its most important property is
\code{.value}: what the user has chosen right now.

\begin{python}[noeval]
athlete_picker = mo.ui.dropdown(
    options=sorted(hr_zoned["name"].unique()),
    value=sorted(hr_zoned["name"].unique())[0],
    label="Athlete",
)
time_window = mo.ui.range_slider(
    start=0, stop=2400, step=60, value=(0, 2400),
    label="Time window (s)", show_value=True,
)
mo.hstack([athlete_picker, time_window])
\end{python}

\begin{itemize}
  \item \code{options} are the choices; \code{value} is the initial
        choice. \code{.unique()} gives every name once, \code{sorted}
        puts them in alphabetical order.
  \item a \code{range\_slider} has two handles; its \code{.value} is a
        tuple, e.g.\ \code{(0, 2400)}.
  \item \code{mo.hstack([...])} puts elements side by side,
        \code{mo.vstack([...])} below each other.
\end{itemize}

In a \emph{different} cell you use the chosen values to filter the data
and make the chart:

\begin{python}[noeval]
selected_hr = hr_zoned.loc[
    (hr_zoned["name"] == athlete_picker.value)
    & hr_zoned["time_s"].between(*time_window.value)
].sort_values("timestamp")

reactive_hr_chart = px.line(selected_hr, x="timestamp", y="heart_rate_bpm",
                            title=f"Heart rate - {athlete_picker.value}")
reactive_hr_chart
\end{python}

\code{.between(a, b)} gives a mask that is \code{True} for values from
\code{a} up to and including \code{b}. The star in
\code{between(*time\_window.value)} \emph{unpacks the tuple}:
\code{between(*(0, 2400))} is the same as \code{between(0, 2400)}.

\begin{keyidea}[Why two cells?]
marimo runs a cell again when a variable it depends on changes. When you
change the dropdown, \code{athlete\_picker.value} changes, and marimo
runs every cell that uses \code{athlete\_picker} again. The cell that
\emph{creates} the element must not read its \code{.value} itself: then
the element would be recreated with every choice. So: \textbf{create a
UI element in one cell, read its value in another}.
\end{keyidea}

\subsection{Selecting in the chart (optional)}

\code{mo.ui.plotly(figure)} goes one step further: when you select
points in the chart with the box or lasso tool, those points are in
\code{.value}. \code{pd.DataFrame(selectable\_scatter.value)} gives them
to you as a table.

% ---------------------------------------------------------------------------
\section{The mini-dashboard}

Today's assignment: a small dashboard for the coach, with \textbf{three
charts} and \textbf{one selector} that at least one chart responds to.
The selector (\code{dashboard\_selector}, a choice between the
positions) is already in the notebook. You make a list
\code{dashboard\_charts} with three figures.

A few tips:

\begin{itemize}
  \item \textbf{Reuse.} The goal is combining, not learning new chart
        types. Anything from this lab can come back.
  \item \textbf{One question per chart.} First write down, for each
        chart, which question of the coach it answers, then choose the
        type (table in section 2).
  \item \textbf{Filter with the selection.} Filter the data with
        \code{dashboard\_selector.value} before you make the chart, as
        with the heart-rate chart above.
  \item \textbf{Finishing.} Every chart has a title and readable axis
        labels with units, and uses \code{POSITION\_COLORS} where
        positions appear.
  \item \textbf{Layout.} \code{mo.vstack} and \code{mo.hstack} arrange
        the selector and the charts neatly below and next to each other.
\end{itemize}

% ---------------------------------------------------------------------------
\section{Summary}

\begin{center}
\rowcolors{2}{tint}{white}
\begin{tabularx}{\textwidth}{@{}L{7.6cm}Y@{}}
\toprule
\textbf{Code} & \textbf{What it does} \\
\midrule
\code{px.histogram(df, x=, nbins=)} & distribution of one variable \\
\code{px.box(df, x=, y=)} & compare distributions per group \\
\code{px.scatter(df, x=, y=)} & relationship between two variables \\
\code{px.line(df.sort\_values(x), x=, y=)} & evolution; sort first \\
\code{px.bar(df, x=, y=)} & compare values; axis from 0 \\
\code{title=}, \code{labels=\{...\}}, \code{hover\_data=[...]} & title, readable axis labels, extra info on hover \\
\code{color=}, \code{color\_discrete\_map=} & category as colour, with fixed colours \\
\code{facet\_col=}, \code{size=} & panels per value, number as point size \\
\code{fig.update\_yaxes(rangemode="tozero")} & let the y-axis start at 0 \\
\code{fig.add\_hrect(y0=, y1=, ...)} & horizontal band, e.g.\ a heart-rate zone \\
\code{mo.ui.dropdown}, \code{mo.ui.range\_slider} & dropdown, slider with two handles \\
\code{element.value} & the user's current choice \\
\code{mo.hstack}, \code{mo.vstack} & put elements side by side or below each other \\
\bottomrule
\end{tabularx}
\end{center}

\subsection*{Common mistakes}

\begin{itemize}
  \item A bar chart whose y-axis does not start at 0.
  \item No title, or no units on the axes.
  \item A line chart on data that is not sorted on the x-axis.
  \item Creating a UI element and reading its \code{.value} in the same
        cell.
  \item Letting every chart choose its own colours, so ``Forward'' is
        blue in one chart and orange in another.
\end{itemize}

\subsection*{Glossary}

\begin{description}[leftmargin=3.2cm, labelwidth=3.0cm, font=\normalfont\bfseries]
  \item[figure] the object a \code{px} function returns; marimo shows it as a chart.
  \item[median] middle value: half of the values are below it, half above.
  \item[quartile] boundary of the lowest or highest 25\,\% of the values.
  \item[outlier] value that lies far from the others.
  \item[facet] small panel with the same axes, one per value of a column.
  \item[UI element] dropdown, slider or other element the user chooses something with.
\end{description}

\end{document}
